% ===== 原文 pp.59--63 =====
% [原文 p.59]
其中 SYM 作用量包含规范耦合与 theta 项，而路径积分测度隐含地计入了规范固定。对于规范化后理论的任何其他可观测值，也有类似的表达式，用未规范化理论的可观测值表示。

接下来的想法是：为矢量多重态加上形变项 (5.11)，但完全不为 tinkertoy 加任何形变项。超对称局域化把矢量多重态的路径积分限制为
\begin{align}
  Z_{\mathrm{T,gauged}} &= \lim_{t\to+\infty}\int[\mathrm{D}A\,\mathrm{D}\lambda\,\mathrm{D}\phi]\,\ee^{-S_{\mathrm{SYM}}-tQV_{\mathrm{vector}}}Z_{\mathrm{T}}[A,\lambda,\phi] \nonumber\\
  &= \int_{QV_{\mathrm{vector}}=0}[\mathrm{D}A\,\mathrm{D}\phi]\,\ee^{-S_{\mathrm{SYM}}[A,\phi]}Z^{\mathrm{vector}}_{\mathrm{one-loop}}[A,\phi]Z_{\mathrm{T}}[A,\phi].
  \tag{5.42}
\end{align}
原则上 $QV_{\mathrm{vector}}=0$ 的轨迹上不能开启费米子 $\lambda$，这就是为什么 $\lambda$ 在最后一行消失了。该轨迹在远离极点处由 (5.18) 描述，在极点处由 (5.19) 描述。也就是说，矢量多重态形变项的零点由常数虚数值 $\phi=a\in\fg$ 以及在极点处的（反）瞬子来刻画。

在我们的应用中，T 分裂为互不耦合的若干扇区 $\mathrm{T}=\prod_L\mathrm{T}_L$。把 $G=\prod_I G_I$ 分解为单因子，并以 $\fh_I$ 记它们的 Cartan 代数，我们得到
\begin{align}
  Z_{\mathrm{T,gauged}} &= \prod_I\left[\sum_{k_I,\bar k_I\ge0}\int_{\fh_I\times\cM_{G_I,k_I}\times\cM_{G_I,\bar k_I}}\dd a_I\,\dd\xi_I\,\dd\bar\xi_I\right]\bigg\{ \nonumber\\
  &\qquad\prod_I\left[|q_I|^{r^2|a_I|^2}q_I^{k_I}\bar q_I^{\bar k_I}Z^{\mathrm{vector}}_{\mathrm{one-loop}}[G_I;a_I,k_I,\xi_I,\bar k_I,\bar\xi_I]\right]\prod_L Z_{\mathrm{T}_L}[a,k,\xi,\bar k,\bar\xi]\bigg\}.
  \tag{5.43}
\end{align}
对 class S 理论，这与在 $C$ 上计算关联函数的方式精确吻合：在 $C$ 的裤子分解的每一根管子上插入一组完备态，并用 $a_I,k_I,\xi_I,\bar k_I,\bar\xi_I$ 标记。确实，我们认出了每根管子 $I$ 上的态求和 $\sum_{k,\bar k}\int_{a,\xi,\bar\xi}$、该态的模之逆，以及对每个三孔球面 $L$ 的其余三点函数 $Z_{\mathrm{T}_L}$。我们知道 $Z_{\mathrm{T}_L}$ 只对 $\mathrm{T}_L$ 带电的那些群所对应的 $a_I,k_I,\xi_I,\bar k_I,\bar\xi_I$ 敏感；在 class S 中，这些正对应于连接到该给定 tinkertoy 的管子。

% [Tutor Note: p.59 (5.43) 与 (5.28)/(5.29) 的逐项对照（核对结构的最快办法）：(a) $\sum_{k_I,\bar k_I\ge0}$ 对应 (5.29) 的 $\sum_{k_I\ge0}$，只是正/反瞬子各有一份；(b) 系数 $q_I^{k_I}\bar q_I^{\bar k_I}$ 对应 (5.29) 的 $\prod_I q_I^{k_I}$ 中同一 $I$ 的那一份；(c) 额外出现的 $|q_I|^{r^2|a_I|^2}$ 正是 (5.28) 中 $Z_{\mathrm{cl}}$ 的整体因子 $|q|^{r^2|a|^2}$ 按管子 $I$ 拆分后的结果（量纲上 $q_I\bar q_I=|q_I|^2$）；(d) $Z^{\mathrm{vector}}_{\mathrm{one-loop}}[G_I;\cdots]$ 对应 (5.24)；(e) $\prod_L Z_{\mathrm{T}_L}$ 对应每个三孔球面的三点函数。所以 (5.43) 就是把 (5.28) 的 $a$ 积分换成“逐管子态求和”的写法，没有新增内容。]

\para{通过在 $S^3_b$ 上局域化实现全纯因子化。} 现在我们转向一种用超对称局域化切割和粘接配分函数的方法。这一方法在 [249, 250] 中被最深入地探讨，那里导出了许多先前观察到的因子化性质（关于该想法的更早应用另见 [91]）。我们这里的目标是：通过把球面沿赤道 $S^3_b\subset S^4_b$ 切成两半，来解释对 $q$ 的全纯依赖。这与我们上面切割黎曼面 $C$ 的做法恰成互补。

作为热身，我们先回顾拉格朗日情形。该情形下因子化的要点是
\begin{itemize}[leftmargin=1.4em,itemsep=2pt,topsep=2pt]
\item 对 $q$ 与 $\bar q$ 的依赖来自极点处的瞬子；
\item 这些瞬子不通过 $Z_{\mathrm{T}_L}$ 相互作用。
\end{itemize}

% [原文 p.60]
准确地说，T 由超多重态组成，而用超多重态形变项作超对称局域化会把 T 的所有场置零，二次涨落给出一个 $a,k,\xi$ 的函数与一个 $a,\bar k,\bar\xi$ 的函数的乘积，正如矢量多重态那样 (5.27)。因此，$q,k,\xi$ 与 $\bar q,\bar k,\bar\xi$ 的依赖在 (5.43) 中完全解耦。把对 $k,\xi$ 的和--积分合成为一个瞬子配分函数，我们就重新得到了所期望的因子化：它是对 $a$ 的积分，被积函数是一个关于 $q$ 的全纯函数乘以一个反全纯函数，正如 (5.28) 所述。

现在我们用局域化的语言重述这一点。4d 场在 $S^3_b$ 上的限制是 3d $\cN=2$ 矢量多重态与超多重态，而完整的 4d 理论也可以等价地描述为 $S^4_b$ 两个半球上的两个（相同的）4d 理论，通过它们的边界条件耦合在一起：例如配分函数可以写为
\begin{equation}
  Z_{S^4_b} = \int[\mathrm{D}\Phi_{3d}]\,\ee^{-S_{3d}[\Phi_{3d}]}Z_{H S^4_b,\mathrm{North}}[\Phi_{3d},q,\bar q]Z_{H S^4_b,\mathrm{South}}[\Phi_{3d},q,\bar q],
  \tag{5.44}
\end{equation}
其中 $\Phi_{3d}$ 统称赤道上 3d $\cN=2$ 理论的所有场，而 $Z_{H S\cdots}$ 表示两个半球上场的路径积分。这个表达式可以看作 3d 理论中（非常复杂的）可观测值 $Z_{H S\cdots}$ 的期望值。

赤道理论随后可以通过给它加上通常用于 $S^3_b$ 上局域化的形变项来局域化。3d 超多重态如 4d 中 (5.15) 那样被局域化到零。3d 矢量多重态被局域化到常数值 $\phi=a\in\fg$。与任何其他可观测值一样，$Z_{H S\cdots}$ 因子在局域化轨迹上取值，于是
\begin{equation}
  Z_{S^4_b} = \int \dd a\, Z_{\mathrm{cl},3d}[a]Z_{\mathrm{one-loop},3d}[a]Z_{H S^4_b,\mathrm{North}}[a,q]Z_{H S^4_b,\mathrm{South}}[a,\bar q].
  \tag{5.45}
\end{equation}
这里半球配分函数是在 Dirichlet 边界条件下计算的，矢量多重态标量的边界值为常数 $a$。3d 中不出现规范耦合依赖是标准的：3d 中的 YM 项是 $Q$-恰当的。更一般地，YM 项与 theta 项在南北极点之外处处是 $Q$-恰当的，而在极点处它们的和或差不能以光滑的方式 $Q$-恰当。这就解释了为什么半球贡献只依赖于 $q$ 或 $\bar q$。

% [Tutor Note: p.60 关于“3d 中 YM 项 $Q$-恰当”的补充论证：3d $\cN=2$ 的 YM 作用量可以写成 $Q$-恰当项（这是 3d 超对称规范理论的一个标准事实，与 4d 中 $\int F\wedge F$ 只在下同调意义下 $Q$-恰当不同）。因此在赤道 $S^3_b$ 上局域化时，$Z_{H S\cdots}$ 对规范耦合 $\tau$ 没有依赖，只剩 $q=\ee^{2\pi\ii\tau}$（即极点瞬子）带来的依赖；这正是 (5.45) 中 $Z_{H S^4_b,\mathrm{North}}$ 只含 $q$、$Z_{H S^4_b,\mathrm{South}}$ 只含 $\bar q$ 的原因。可核对之处：若把 $\tau\to\tau+1$，四维配分函数不变（$q\to q$），而 $q\to\ee^{2\pi\ii\tau}$ 的全纯因子在 $\tau\to\tau+1$ 下也不变，二者一致。]

这个沿赤道切割理论的想法可以推广到拉格朗日理论之外。即使 $Z_{\mathrm{T}_L}$ 可能没有路径积分描述，QFT 的一般公理仍然适用，人们可以沿超曲面 $S^3_b\subset S^4_b$ 插入一组完备态。于是配分函数可看作由两个半球上的路径积分所制备的两个态的内积，
\begin{equation}
  Z_{S^4_b} = \sum_{|\Phi\rangle\in\cH[S^3_b]}\langle H S^4_b\ \mathrm{North}\,|\,\Phi\rangle\langle\Phi\,|\,H S^4_b\ \mathrm{South}\rangle.
  \tag{5.46}
\end{equation}
接着，该框架的 $Q$-不变性意味着求和只限于 $Q$-不变态，
\begin{equation}
  Z_{S^4_b} = \sum_{|a\rangle\in\cH[S^3_b],\,Q|a\rangle=0}\langle H S^4_b\ \mathrm{North}\,|\,a\rangle\langle a\,|\,H S^4_b\ \mathrm{South}\rangle,
  \tag{5.47}
\end{equation}
% [原文 p.61]
而每个半球上 YM $\pm$ theta 项是 $Q$-恰当的，这表明两个因子正是所要求的 $q$ 的（反）全纯函数。

我们注意到，对 $a\in\fh$ 的 Coulomb 分支积分一般而言被替换为对 4d 理论 $S^3_b$ Hilbert 空间中 $Q$-不变态的求和/积分。这通常是一个更大的空间，特别是在存在非拉格朗日 tinkertoy 时。

% [Tutor Note: p.60--61 (5.46)--(5.47) 是“用完备态代替积分”的关键一步：原来的 $\int\dd a$ 只遍历 $Q$-不变态的连续谱（Coulomb 分支），而 (5.47) 的求和遍历整个 $S^3_b$ Hilbert 空间的 $Q$-不变子空间，后者在存在非拉格朗日 tinkertoy 时严格更大。这个替换之所以合法，是因为 $Q$-恰当形变项的局域化只要求“在局域化轨迹上取值”，而轨迹由 $Q$-不变态构成；(5.45) 只是 (5.47) 在拉格朗日情形（此时 $Q$-不变态恰由 $a$ 参数化）的特例。]

\section{SU(2) 箭图的 AGT 对应}

AGT 关系 (5.1) 把 6d $(2,0)$ 理论 $\cX(\fsu(2))$ 的两个不同维数约化的可观测值联系起来。我们在第 4 节已解释：把 $\cX(\fsu(2))$ 在黎曼面 $\cC=\overline{\C}\setminus\{z_1,\dots,z_n\}$ 上约化，会得到一个 4d $\cN=2$ $\fsu(2)$ Sicilian\footnote{之所以这样命名，是因为 trinion（三孔球面）与西西里旗帜上醒目的三曲腿图（triskelion）形状相似。}箭图规范理论，它依赖于 $\cC$，其规范耦合与 $\cC$ 的复结构相联系。同样地，我们预期把 $\cX(\fsu(2))$ 在 $S^4_b$ 上约化应得到一个耦合常数为 $b$ 的 2d 理论，而插入在 $\cC$ 的穿刺点 $z_i$ 处的 $\cX(\fsu(2))$ 余维 2 缺陷应成为 2d 中的局域算子。由于这个 2d 维数约化略显技术性，我们把它推迟到 6.3 小节，并在那里简要说明为什么应当预期 6d 可观测值在 4d 侧与 2d 侧都可计算。

在那之前，我们在 6.2 小节中以更符合历史的方式来判定相关的 2d 理论：它的关联函数应当重现 class S 理论的 $S^4_b$ 配分函数。这些配分函数只依赖于 $\cC$ 的复结构，因而只依赖于 $\cC$ 上度规的共形类。这意味着我们所寻求的 2d 理论应当是一个 CFT，而它最终证明就是 Liouville 共形场论。因此，我们首先在 6.1 小节回顾 Liouville 共形场论与 2d CFT 的基础。我们在表 3 中总结了该对应关系的若干方面。

\begin{table}[htbp]
\centering
\small
\renewcommand{\arraystretch}{1.2}
\begin{tabular}{@{}cll@{}}
\toprule
 & \textbf{规范理论} & \textbf{Toda/Liouville 共形场论} \\
\midrule
 & $S^4_b$ 上的配分函数 & $\cC$ 上的 Toda 关联函数 \\
\midrule
\multirow{4}{*}{\rotatebox{90}{参数}} & 压扁参数 $b=\sqrt{\epsilon_1/\epsilon_2}$ & Toda 耦合常数 $b$ \\
 & 4d 规范耦合 & $\cC$ 的复结构 \\
 & 对偶规范理论描述 & $\cC$ 的裤子分解 \\
 & 质量为 $m$ 的满驯服穿刺 & 顶点算子 $\hat V_{Q+rm}$ \\
\midrule
\multirow{5}{*}{\rotatebox{90}{构建单元}} & Coulomb 分支参数 $a$ & 内部动量 $\alpha=Q+ra$ \\
 & 规范单圈行列式 & 两点函数之逆 \\
 & 物质单圈行列式 & 三点结构常数 \\
 & 经典作用量贡献 & 共形块领头项 \\
 & 瞬子配分函数 & 共形块幂级数 \\
\bottomrule
\end{tabular}
\tabtitle{3}{AGT 的要点。在 Liouville 共形场论的约定下，把 $Q$ 换成 $Q/2$。}
\end{table}

% [Tutor Note: p.61 表 3 是原文中真实存在的第三个表（STYLE_GUIDE §7 的原清单只列了表 1、表 2 与缩写表，此处已按实际补上，请 Lead 一并核对 main.tex 是否有表清单需要更新）。表中两点值得注意：(1) 表题自注“在 Liouville CFT 约定下把 $Q$ 换成 $Q/2$”，是因为标准 Liouville 约定取 $c=1+6Q^2$、$Q=b+1/b$，顶点算子 $V_\alpha$ 的共形权为 $\alpha(Q-\alpha)$；把表中的 $\alpha=Q+ra$ 换成 $\alpha=Q/2+ra$ 后 $h=\alpha(Q-\alpha)=Q^2/4-r^2a^2$，对虚的 $a$ 为正，与 4d 侧 $r^2|a|^2$ 的符号一致。(2) “满驯服穿刺”（full tame puncture）的质量 $m$ 对应顶点算子 $\hat V_{Q+rm}$，$\hat V$ 是本节后面出现的 AGT 友好归一化（p.63 也写作 $\hat V_\alpha$）。]

% [原文 p.62]
\subsection{二维 CFT 与 Liouville 共形场论}

关于 2d CFT（特别是 Liouville 共形场论）的入门，可参见 [128, 129] 等综述（及其中的参考文献）。青年研究者可积性与工作坊（Young Researchers Integrability School and Workshop，YRISW）课程 [251] 也给出了 2d CFT 及其对称代数（手征代数）的简要介绍，而这些对称代数还奇特地独立出现在 4d $\cN=2$ 理论的受保护子扇区中。

\para{Virasoro 代数。} 与其他维数不同，2d 中的共形对称代数是两个无穷维代数之积 $\mathrm{Vir}\times\overline{\mathrm{Vir}}$。Virasoro 代数 $\mathrm{Vir}$ 由 $L_n$（$n\in\Z$）与一个中心元 $C$ 张成，并满足
\begin{equation}
  [L_m,L_n] = (m-n)L_{m+n}+\frac{1}{12}\delta_{m+n=0}(m^3-m)C.
  \tag{6.1}
\end{equation}
这个代数，以及由 $\bar L_n$（$n\in\Z$）与 $\bar C$ 张成的另一份拷贝，作用在圆上理论的 Hilbert 空间上。对任何给定的 2d CFT，共形代数的元素 $C,\bar C$ 都按常数相乘作用，这些常数称为中心荷，记作 $c,\bar c$ 或 $c_L,c_R$。

在 CFT 中，径向量子化把这些态等同于它们在径向上无穷远的过去的极限，而那是径向量子化中心处的一个局域算子 $\cO$。在此态--算子对应下，与 $\cO$ 对应的态记为 $|\cO\rangle$。在这一对应下，$\mathrm{Vir}\times\overline{\mathrm{Vir}}$ 对态的作用翻译为通过交换子对局域算子的作用。

\begin{exercise}{6.1}
利用 (6.1) 验证 $[L_m,[L_n,L_p]]+[L_n,[L_p,L_m]]+[L_p,[L_m,L_n]]=0$。这就是 Jacobi 恒等式，它对 Lie 代数 $\mathrm{Vir}$ 的自洽性是必需的。验证因子 $m^3-m$ 由 Jacobi 恒等式在“把 $L_n$ 与 $C$ 混合、并对 $C$ 标度”的意义下被唯一确定。

验证 $\tilde C=kC$ 与 $\tilde L_n=L_{kn}/k+\delta_{n=0}(k-1/k)C/24$ 对任意 $k\in\Z\setminus\{0\}$ 都满足 Virasoro 对易关系，从而 Virasoro 代数包含无穷多个 Virasoro 子代数。反过来，Virasoro 代数可以作为具有分数模的更大 Virasoro 代数的整数模嵌入其中，这在对称积轨形中很有用 [252]。
\end{exercise}

% [Tutor Note: p.62 关于 Exercise 6.1 的核对（可逐项验算）：设 $\tilde C=kC$、$\tilde L_n=L_{kn}/k+\delta_{n=0}(k-1/k)C/24$。则 $[\tilde L_m,\tilde L_n]=\frac{1}{k^2}[L_{km},L_{kn}]=\frac{m-n}{k}L_{k(m+n)}+\frac{1}{12}\delta_{m+n=0}(m^3k-m/k)C$。把第一项写成 $(m-n)\big(\tilde L_{m+n}-\delta_{m+n=0}(k-1/k)C/24\big)$，再与所需的 $\frac{1}{12}\delta_{m+n=0}(m^3-m)\tilde C$ 比较，利用 $m+n=0$ 时 $m-n=2m$，可得 $C$ 的系数为 $\frac{1}{12}mk(m^2-1)$，恰好等于 $\frac{k}{12}(m^3-m)$------即 $\tilde C=kC$ 正是中心荷。这同时说明 $m^3-m$ 的系数 $1/12$ 与 $C$ 的标度方式一致，是 Jacobi 恒等式所固定的。]

\para{共形维数。} 绕径向量子化中心的伸缩与旋转由 $L_0\pm\bar L_0$ 生成，因此局域算子的维数 $\Delta=h_\cO+\bar h_\cO$ 与自旋 $h_\cO-\bar h_\cO\in\frac{1}{2}\Z$ 由 $L_0$ 与 $\bar L_0$ 的作用给出：
\begin{equation}
  [L_0,\cO] = h_\cO\cO,\qquad [\bar L_0,\cO]=\bar h_\cO\cO.
  \tag{6.2}
\end{equation}
尽管记号如此，共形维数 $h_\cO$ 与 $\bar h_\cO$ 是互相独立的数。在幺正 CFT 中它们都是实的且非负。我们发现主要用态来处理更方便，此时 $\Delta_\cO$ 与 $h_\cO-\bar h_\cO$ 分别是态的能量与动量。

有趣的是，对易子 $[L_0,L_n]=-nL_n$ 意味着：若 $|\cO\rangle$ 的共形维数为 $(h_\cO,\bar h_\cO)$，则 $L_n|\cO\rangle$ 的共形维数为 $(h_\cO-n,\bar h_\cO)$。因此 $L_n$（$n\ge1$）称为降算子，$L_n$（$n\le-1$）称为升算子。对 $\bar L_n$ 同理。

% [原文 p.63]
\para{主算子。} 给定理论的 Hilbert 空间组织成共形族，即 $\mathrm{Vir}\times\overline{\mathrm{Vir}}$ 的表示。在幺正 CFT 中态的能量非负，因此每个共形族都有一个极小维数的态 $|V\rangle$。这样的态被所有降算子 $L_n,\bar L_n$（$n\ge1$）湮灭，称为主态（相应的算子 $V$ 称为主算子）：
\begin{equation}
  L_n|V\rangle = \bar L_n|V\rangle = 0,\qquad n\ge1.
  \tag{6.3}
\end{equation}
从这个维数为 $(h,\bar h)$ 的态出发，可以用 $L_{-n}$ 与 $\bar L_{-n}$（$n\ge1$）作用，构造出维数更高的态的塔。利用 Virasoro 对易关系 (6.1)，这些升算子可以被排序，而后代态所张成的集合由如下态给出
\begin{equation}
  \mathbb{L}_{-Y}\bar{\mathbb{L}}_{-\bar Y}|V\rangle := L_{-m_1}\dots L_{-m_k}\bar L_{-n_1}\dots\bar L_{-n_l}|V\rangle,
  \tag{6.4}
\end{equation}
其中 $Y,\bar Y$ 是两个杨图，$m_1\ge m_2\ge\cdots\ge m_k\ge1$ 是 $Y$ 的各行长度（依次），而 $n_1\ge\cdots\ge1$ 类似地是 $\bar Y$ 的各行长度。(6.4) 的共形维数为 $(h+m_1+\cdots+m_k,\ \bar h+n_1+\cdots+n_l)$。这些态称为 $|V\rangle$ 的后代（descendants）。事实上，$\mathrm{Vir}\times\overline{\mathrm{Vir}}$ 的整个表示（作为向量空间）都由 (6.4) 张成，而一般说来这些态线性无关。

\begin{exercise}{6.2}
1. 设 $|V\rangle$ 是主态，对 $m,n\in\Z$，把 $L_mL_n|V\rangle$ 重写为指标已正确排序的 (6.4) 型项之线性组合。

2. 验证：对任意 $Y,\bar Y$，用任意 $L_n$ 或 $\bar L_n$ 作用于 $\mathbb{L}_{-Y}\bar{\mathbb{L}}_{-\bar Y}|V\rangle$ 都得到这类后代态的线性组合。
\end{exercise}

\para{两点与三点函数。} 共形对称性约束局域算子的关联函数。记 $z_{ij}=z_i-z_j$，则维数为 $(h_i,\bar h_i)$ 的主算子 $V_i$ 的两点与三点函数取如下形式
\begin{align}
  \langle V_1(z_1,\bar z_1)V_2(z_2,\bar z_2)\rangle &= g_{12}\,\delta_{h_1=h_2}\delta_{\bar h_1=\bar h_2}\,z_{12}^{-2h_1}\bar z_{12}^{-2\bar h_1}, \nonumber\\
  \langle V_1(z_1,\bar z_1)V_2(z_2,\bar z_2)V_3(z_3,\bar z_3)\rangle &= C_{123}\,z_{23}^{h_1-h_2-h_3}z_{31}^{h_2-h_3-h_1}z_{12}^{h_3-h_1-h_2} \nonumber\\
  &\qquad\times\bar z_{23}^{\bar h_1-\bar h_2-\bar h_3}\bar z_{31}^{\bar h_2-\bar h_3-\bar h_1}\bar z_{12}^{\bar h_3-\bar h_1-\bar h_2},
  \tag{6.5}
\end{align}
其中 $g_{12},C_{123}$ 是依赖于所涉主算子的常数。后代算子可以写成作用在主算子上的 Virasoro 生成元 $L_n,\bar L_n$，而这些生成元按各种微分算子作用在 (6.5) 上。

若 CFT 对每个共形维数 $(h,\bar h)$ 都只有一个主算子（Liouville 共形场论正是如此），则 $g_{12}$ 可以被吸收进该算子的归一化中。尽管有这一可能，主算子 $V_\alpha$ 的标准归一化给出的 $g_{12}$ 仍是非平凡的，而我们下面给出的 AGT 友好归一化 $\hat V_\alpha$ 也如此。归一化一经选定，理论就由其主算子谱（即它们的共形维数）以及三点函数 $C_{123}$ 所刻画，这将在下文中说明。

% [Tutor Note: p.63 (6.5) 的指标核对：两点函数里 $z_{12}^{-2h_1}$ 与 $\bar z_{12}^{-2\bar h_1}$ 的指数必须分别是 $h_1,\bar h_1$ 的两倍，这来自 $L_0$ 作用两次（即 $z^{-h}$ 每因子）；三点函数的指数满足“对每个 $i$，$\sum_j(\text{指数}_{ij})=-2h_i$”：例如对 $i=1$，$z_{12}$ 的指数 $h_3-h_1-h_2$ 加 $z_{31}$ 的指数 $h_2-h_3-h_1$ 得 $-2h_1$ ✓。这正是共形不变性（$z_i\to az_i+b$ 与反演）唯一确定三点函数形式的标准论证，可用来快速排查抄写错误。另注意 (6.5) 的三行在原图中是一个编号 (6.5)，其编号位于三点函数的第二行右侧。]
